\section{Formal Foundations and Taxonomy}
\label{sec:taxonomy}

\subsection{Basic Mathematical Constructs}
\label{subsec:basics}

\subsubsection{Encoding Tree of a Finite Set}

Let $A$ be a finite set representing the entities in a complex system (e.g.,
nodes in a network or data points in a dataset).
A natural model for the structure of such a system is a \emph{hierarchical
abstraction} in which elements are recursively grouped into nested subsets.

\begin{definition}[Encoding Tree]
\label{def:encoding-tree}
Given a finite set $A$, an \emph{encoding tree} of $A$ is a rooted tree $T$
such that:
(1) the root $\lambda$ corresponds to the entire set $T_\lambda = A$;
(2) for every non-leaf node $\alpha\in T$ with children
    $\beta_0,\ldots,\beta_l$ the corresponding subsets form a partition
    $T_\alpha = \bigsqcup_{j=0}^l T_{\beta_j}$;
(3) each leaf $\gamma\in T$ corresponds to a singleton $T_\gamma = \{a\}$.
\end{definition}

The \emph{dimension} of an encoding tree is its depth.
The extreme case---an encoding tree of depth~1 where the root is directly
partitioned into singleton leaves---yields the \textbf{1-dimensional SE},
which reduces to Shannon entropy.

\subsubsection{Structural Entropy}

Consider a non-negative, irreducible matrix
$A\in\mathbb{R}_{\ge 0}^{n\times n}$ representing an information system over
$n$ elements indexed by $V=\{1,\ldots,n\}$.
Let $\boldsymbol\pi=(\pi_1,\ldots,\pi_n)$ be the stationary distribution of
the Markov chain derived from $A$ by row normalization, with transition
probability $p_{xy}=\pi_x\,b_{xy}$ where $b_{xy}=a_{xy}/\sum_j a_{xj}$.

For a subset $X\subseteq V$ define:
\[
  V_X = \sum_{x\in X}\pi_x, \qquad
  p_X = \sum_{y\notin X}\sum_{x\in X}p_{yx}, \qquad
  q_X = \sum_{x\in X}\sum_{y\notin X}p_{xy}.
\]

\begin{definition}[Structural Entropy under a Tree]
\label{def:se-tree}
The \emph{structural entropy} of $A$ under encoding tree $T$ is
\[
  \mathcal{H}^T(A)
  = -\sum_{\substack{\alpha\in T\\\alpha^-\neq\lambda}}
    p_\alpha\log_2\frac{V_\alpha}{V_{\alpha^-}},
\]
where $\alpha^-$ denotes the parent of $\alpha$ in $T$.
\end{definition}

\begin{definition}[Structural Entropy of a System]
\label{def:se-system}
The \emph{structural entropy} of an information system is
\[
  \mathcal{H}(A) = \min_T \mathcal{H}^T(A).
\]
\end{definition}

The minimizing tree $T^*$ encodes the system's most efficient
information-theoretic abstraction---its natural hierarchical community
structure~\cite{Li2016}.

\subsection{Relationship to Other Graph-Complexity Measures}
\label{subsec:relations}

Three major graph-complexity measures are related to SE:

\paragraph{Modularity~\cite{newman2004findingmodularity}.}
Modularity $Q$ measures the fraction of edges within communities minus the
expected fraction under a null model.
While modularity seeks to maximize within-community density relative to a
random graph, SE minimization seeks to \emph{minimize informational
uncertainty} within a hierarchical partition.
Modularity is resolution-limited (it cannot detect communities smaller than a
scale-dependent threshold); SE is not subject to the same limitation.

\paragraph{Map Equation~\cite{rosvall2008maps}.}
Infomap minimizes the expected description length of a random walk under a
two-level codebook.
SE generalization replaces the two-level restriction with an arbitrary tree,
and replaces flow coding with a volume-and-cut measure.
Formally, \cite{liu2022similaritygap} proved that for any undirected graph
$\mathcal{H}(A) \le L^*(A) \le \mathcal{H}(A) + \gamma$, where $L^*$ is the
optimal map-equation description length and $\gamma$ depends only on the
graph's degree sequence.

\paragraph{Von Neumann Entropy (VNE).}
VNE is the entropy of the normalized graph Laplacian's eigenvalue
distribution.
\cite{liu2022similaritygap} established that
$\mathcal{H}(A) - \mathrm{VNE}(A) \in [0,\gamma]$, providing the first
formal spectral interpretation of SE.

\subsection{Taxonomy of SE Research}
\label{subsec:taxonomy}

Figure~\ref{fig:taxonomy} provides a high-level taxonomy organizing SE
research into three pillars---\emph{Theory}, \emph{Learning Methods}, and
\emph{Cross-Domain Applications}---each containing major sub-themes.
We follow this taxonomy throughout the survey.

\begin{figure}[ht]
\centering
\fbox{\parbox{0.9\textwidth}{%
\textbf{Structural Entropy Research}
\begin{itemize}
  \item \textbf{Theory}
    \begin{itemize}
      \item Foundational works: SE~\cite{Li2016}, ISPAI~\cite{li2024book}
      \item Theoretical analysis: spectral gap~\cite{liu2022similaritygap},
            incremental~\cite{yang2024incremental},
            togetherness~\cite{ljm2021improvingTogetherness}
      \item Extensions: directed~\cite{wxb2024rwedirected,yaopan2019flowDirected},
            multi-relational~\cite{cao2024multi},
            source localization~\cite{wxb2023findingThesource},
            probabilistic coding~\cite{huang2024structuralProbCoding}
    \end{itemize}
  \item \textbf{Computational Algorithms}
    \begin{itemize}
      \item Heuristic: deDoc~\cite{li2018dedoc}, SEAT~\cite{chen2022pyseat},
            HCSE~\cite{pan2021information}
      \item Constrained: SuperTAD~\cite{zhang2021supertad},
            deDoc2~\cite{li2023dedoc2},
            SuperTAD2~\cite{ling2024supertad2},
            SSE~\cite{zeng2024semiSSE},
            SSSE~\cite{zeng2024scalableSSSE},
            CoDeSEG~\cite{xian2025communitySEGame}
      \item Differentiable: LSENet~\cite{sun2024lsenet},
            DeSE~\cite{zhang2025dese}
    \end{itemize}
  \item \textbf{Learning Methods}
    \begin{itemize}
      \item Graph pooling: SEP~\cite{wu2022pooling},
            Hi-PART~\cite{wxb2024hi-part}, EDEN~\cite{li2025eden}
      \item Structure augmentation: SE-GSL~\cite{zou2023segsl},
            SEGA~\cite{wu2023sega},
            USER~\cite{wang2023user},
            SeGSL~\cite{duan2024structuralGSL}
      \item Kernels \& dimensions: HAGK~\cite{yang2024graphkernal},
            MGEDE~\cite{yang2023MGEDE}
      \item Reinforcement learning: SIRD~\cite{zeng2023effective},
            SISA~\cite{zeng2023hierarchicalSISA},
            SI2E~\cite{zeng2024effectiveSI2E},
            COLLAB~\cite{su2025emergencecollab}
    \end{itemize}
  \item \textbf{Applications}
    \begin{itemize}
      \item Bioinformatics: cancer~\cite{li2016cancer},
            TAD detection~\cite{zhang2021supertad,li2023dedoc2},
            single-cell~\cite{chen2022pyseat,li2021detoki}
      \item Transport \& geoscience:
            traffic~\cite{zou2024multispans},
            scientific knowledge~\cite{wxb2022scientificxray,wang2023quantifyingKnowledge}
      \item Social networks: bot detection~\cite{peng2024zjyun,yang2024sebot},
            event detection~\cite{cao2024socialevent,yang2024eventdetect,yu2024hyperSED},
            safety~\cite{zeng2025si2af}
      \item Pattern recognition: image~\cite{zeng2023lesion,xie2025superpixel},
            text~\cite{zhu2023HiTIN,huang2025MASGCN},
            speech~\cite{wang2024secodecSpeech,wang2025voice}
    \end{itemize}
\end{itemize}
}}
\caption{Taxonomy of Structural Entropy research (2016--2026).}
\label{fig:taxonomy}
\end{figure}

